\section{Набор данных}

Обучение модели выполнено на синтетическом наборе данных, сгенерированном под обучение разработанной архитектуры модели. Набор данных включает `Train', `Val' и `Test'-коллекций csv-таблиц частотных характеристик и json-файлов с координатами полюсов/нулей.

Тренировочный набор данных состоит из 72 случайных комбинаций полюсов/нулей обобщенной модели линейного динамического звена:
\begin{equation*}
    H(s) = G \frac{\prod\limits_{m=1}^{M} \left(1 + \frac{s}{\omega_{\text{z}(m)}} \right)}{s^K \prod\limits_{n=1}^{N} \left(1 + \frac{s}{\omega_{\text{p}(n)}}\right)} e^{-s t_\text{delay}},
\end{equation*}
где
\begin{itemize}
    \item ${G}$ -- пропорциональный коэффициент, ${G \in \mathbb{R}}$;
    \item ${t_\text{delay}}$ -- задержка времени, ${t_\text{delay}>0}$;
    \item ${K}$ -- количество интеграторов, ${K \in \mathbb{N}}$;
    \item ${\omega_{\text{z}(m)}}$ -- частоты нулей передаточной функции ${H(s)}$ в количестве ${M}$, ${\omega_{\text{z}(m)} \in \mathbb{R}}$;
    \item ${\omega_{\text{p}(n)}}$ -- частоты полюсов передаточной функции ${H(s)}$ в количестве ${N}$, ${\omega_{\text{p}(n)} \in \mathbb{R}}$.
\end{itemize}
Каждой комбинации соответствует 500 значений, отличающихся диапазонами частот и частотами полюсов/нулей. Параметры ${G}$, ${t_\text{delay}}$ заданы константами для последующего изменения в ходе обучения. Общий размер тренировочного набора данных -- 72x500=36~000 примеров. Размеры валидационного и тестового наборов -- 14~400 и 7~200 примеров, соответственно.

В тестовом наборе данных параметры ${G}$, ${t_\text{delay}}$ сгенерированы случайным образом; частотные характеристики приближены к реальным данным путем введения шума с нормальным распределением.

Полное описание генератора набора данных приведено в репозитории \cite{gorbunov2026zerospolesmasks}.